\documentclass{beamer}
\usepackage[utf8]{inputenc}
\usepackage[T1]{fontenc}
\usepackage[spanish]{babel}
\usepackage{graphicx}

\setcounter{secnumdepth}{2} % Profundidad de secciones y subsecciones
\makeatletter

\renewcommand\thesection{\arabic{section}}              % Ej: "1"   en lugar de "Section 1"
\renewcommand\thesubsection{\thesection.\arabic{subsection}}  % Ej: "1.1"
\makeatother

 %---------- OPCIÓN 3: SECCIÓN.SUBSECCIÓN.FRAME ----------
 %(sólo si manejas secciones/subsecciones numeradas)
\setbeamertemplate{frametitle}{%
  \begin{beamercolorbox}[wd=\paperwidth,ht=2.25ex,dp=1ex]{frametitle}%
    \quad \thesubsection\quad \insertframetitle
  \end{beamercolorbox}%
}

 
% ==== SECCIONADO AUTOMÁTICO EN DIAPOSITIVAS ====
% Muestra una diapositiva con el título de la sección al comenzar cada \section
\AtBeginSection[]{
    \begin{frame}[plain]
        \centering
        \vfill
        {\usebeamerfont{section title}\Huge \insertsectionhead}
        \vfill
    \end{frame}
}

% Muestra una diapositiva con el título de la subsección al comenzar cada \subsection
\AtBeginSubsection[]{
    \begin{frame}[plain]
        \centering
        \vfill
        {\usebeamerfont{subsection title}\Large \insertsubsectionhead}
        \vfill
    \end{frame}
}
% ==============================================

\usetheme{Madrid} % Puedes cambiar el tema según prefieras

\title{Introducción a los Sistemas Operativos}
\author{Víctor de Juan}
\date{Curso 2024-2025}

\setcounter{secnumdepth}{2} % Profundidad de secciones y subsecciones

\begin{document}

%-------------
\begin{frame}
    \titlepage
\end{frame}

% ------------------
% PRIMERA SECCIÓN (antes del índice)
% ------------------

\begin{frame}{Índice General}
    \tableofcontents
\end{frame}

% ------------------
% SEGUNDA SECCIÓN
% ------------------

%----------------------------------------------------------------------------------------
% 1. INTRODUCCIÓN
%----------------------------------------------------------------------------------------
\section{Introducción}
\subsection{Sistema Operativo}
\begin{frame}{¿Qué es un Sistema Operativo?}
    \begin{itemize}
        \item Un sistema operativo (SO) es un software fundamental que actúa como intermediario entre el usuario y el hardware.
        \item Gestiona y coordina los recursos del sistema (CPU, memoria, dispositivos) de forma eficiente.
        \item Simplifica la interacción del usuario y de los programas con los componentes físicos de la computadora.
        \item Ofrece un conjunto de servicios y abstracciones (archivos, procesos, etc.) que facilitan el desarrollo y la ejecución de aplicaciones.
    \end{itemize}
    \vfill
    \centering
    \textbf{¿Un microondas tiene sistema operativo? ¿Y una lavadora?}
\end{frame}

%----------------------------------------------------------------------------------------
% FRAME: FIRMWARE, SO Y DRIVER
%----------------------------------------------------------------------------------------
\subsection{Firmware, SO y Drivers}
\begin{frame}{Firmware, SO y Driver: Diferencias (1/3)}
    \begin{itemize}
        \item \textbf{Firmware}:
            \begin{itemize}
                \item Software de muy bajo nivel, grabado en memoria no volátil (ROM, EEPROM, Flash).
                \item Controla las funciones básicas de un dispositivo sin necesidad de un sistema operativo.
                \item Se actualiza con poca frecuencia y está muy ligado al hardware.
                \item \textbf{Ejemplos reales}:
                    \begin{itemize}
                        \item BIOS/UEFI de una PC, que permite el arranque del sistema.
                        \item Firmware de un router, que maneja la conectividad sin necesidad de un SO adicional.
                        \item Firmware de una lavadora inteligente, que gestiona los ciclos de lavado.
                    \end{itemize}
            \end{itemize}
    \end{itemize}
\end{frame}

\begin{frame}{Firmware, SO y Driver: Diferencias (2/3)}
    \begin{itemize}
        \item \textbf{Sistema Operativo (SO)}:
            \begin{itemize}
                \item Software que gestiona recursos (CPU, memoria, dispositivos) y ofrece servicios a otras aplicaciones.
                \item Se ejecuta normalmente sobre el hardware y facilita la interacción con el usuario y otros programas.
                \item Puede manejar múltiples procesos y \textbf{coordinar} el uso del hardware.
                \item \textbf{Ejemplos reales}:
                    \begin{itemize}
                        \item \textbf{Windows 11}: Uso en PC y entornos empresariales.
                        \item \textbf{Linux (Ubuntu, Debian, Red Hat, Fedora, Arch-Linux...)}: Presente en servidores, IoT y PC.
                        \item \textbf{Android (basado en Linux)}: Utilizado en smartphones y televisores inteligentes.
                        \item \textbf{FreeRTOS}: Utilizado en sistemas embebidos como electrodomésticos y dispositivos IoT.
                    \end{itemize}
            \end{itemize}
    \end{itemize}
\end{frame}

\begin{frame}{Firmware, SO y Driver: Diferencias (3/3)}
    \begin{itemize}
        \item \textbf{Driver}:
            \begin{itemize}
                \item Programa o módulo que actúa como traductor entre el SO y un hardware específico.
                \item Cada dispositivo (impresora, tarjeta de red, GPU, etc.) necesita su driver para funcionar correctamente.
                \item Gestiona las operaciones de E/S y la comunicación a bajo nivel con el dispositivo.
                \item Puede estar incluido en el sistema operativo o requerir instalación manual.
                \item \textbf{Ejemplos reales}:
                    \begin{itemize}
                        \item \textbf{Driver de GPU (NVIDIA, AMD, Intel)}: Permite la aceleración 3D y el procesamiento gráfico avanzado.
                        \item \textbf{Driver de impresora HP o Epson}: Traducen comandos de impresión en instrucciones específicas para el hardware.
                        \item \textbf{Driver de audio (Realtek, Sound Blaster)}: Controla la salida y entrada de sonido en el sistema.
                        \item \textbf{Driver genérico HID (Human Interface Device)}: Permite el uso de teclados y ratones sin necesidad de instalación manual.
                    \end{itemize}
            \end{itemize}
    \end{itemize}
\end{frame}

%----------------------------------------------------------------------------------------
% SECCIÓN: EL PROCESADOR Y LOS DRIVERS
%----------------------------------------------------------------------------------------
\section{El Procesador y los Drivers}

%----------------------------------------------------------------------------------------
% SUBSECCIÓN: VISIÓN GENERAL
%----------------------------------------------------------------------------------------
\subsection{Visión General}

\begin{frame}{¿Dónde se ejecutan los drivers?}
    \begin{itemize}
        \item Un driver es un software que se ejecuta en la \textbf{CPU del sistema} (procesador principal).
        \item Se encarga de \textbf{comunicarse con el dispositivo de hardware}, traduciendo las peticiones del sistema operativo.
        \item La complejidad del driver depende del \textbf{tipo de dispositivo}: desde periféricos simples (ratón, teclado) hasta sistemas avanzados (GPU).
    \end{itemize}
\end{frame}

%----------------------------------------------------------------------------------------
% SUBSECCIÓN: DRIVERS SIMPLES
%----------------------------------------------------------------------------------------
\subsection{Drivers Simples}

\begin{frame}{Ejemplo: Ratón o Teclado}
    \begin{itemize}
        \item Dispositivos de entrada con \textbf{baja complejidad}.
        \item El driver genérico (HID) se encarga de interpretar movimientos y clics.
        \item \textbf{Flujo de funcionamiento}:
            \begin{enumerate}
                \item El usuario mueve el ratón (o pulsa una tecla).
                \item El ratón/teclado envía señales (USB, Bluetooth, etc.) con coordenadas o pulsaciones.
                \item La \textbf{CPU} procesa la señal mediante el \textbf{driver HID}.
                \item El \textbf{SO} actualiza la posición del cursor o muestra la pulsación.
            \end{enumerate}
        \item Bajo consumo de recursos y respuestas rápidas.
    \end{itemize}
\end{frame}

%----------------------------------------------------------------------------------------
% SUBSECCIÓN: DRIVERS INTERMEDIOS
%----------------------------------------------------------------------------------------
\subsection{Drivers Intermedios}

\begin{frame}{Ejemplo: Tarjeta de Sonido o Impresora}
    \begin{itemize}
        \item La CPU realiza más trabajo que en un ratón/teclado.
        \item \textbf{Tarjeta de sonido}:
            \begin{itemize}
                \item Recibe datos de audio (MP3, WAV, etc.) de una aplicación.
                \item El \textbf{driver} convierte estos datos en señales que la tarjeta pueda procesar.
                \item La tarjeta de sonido reproduce el audio mediante su DAC (Digital-to-Analog Converter).
                \item Algunas tarjetas tienen \textbf{DSP} (Digital Signal Processor), que asume parte del procesamiento.
            \end{itemize}
        \item \textbf{Impresora}:
            \begin{itemize}
                \item El driver traduce los documentos en un formato comprensible para el hardware.
                \item Controla detalles como resolución, color, y manejo de páginas.
            \end{itemize}
        \item Este tipo de drivers implica una \textbf{carga media} para la CPU.
    \end{itemize}
\end{frame}

%----------------------------------------------------------------------------------------
% SUBSECCIÓN: DRIVERS COMPLEJOS
%----------------------------------------------------------------------------------------
\subsection{Drivers Complejos}

\begin{frame}{Ejemplo: Tarjeta Gráfica (GPU)}
    \begin{itemize}
        \item Una GPU moderna tiene su propio procesador especializado y memoria dedicada.
        \item El \textbf{driver gráfico} es muy complejo y gestiona:
            \begin{itemize}
                \item Renderizado 3D, shaders, cálculo de física, codificación/decodificación de video.
                \item Coordinación entre CPU y GPU para dividir cargas de trabajo.
            \end{itemize}
        \item \textbf{Flujo en un videojuego}:
            \begin{enumerate}
                \item La CPU recibe instrucciones para generar gráficos.
                \item El \textbf{driver} envía comandos a la GPU (renderizado, texturas, efectos).
                \item La GPU procesa en \textbf{paralelo}, genera las imágenes y las envía al monitor.
                \item El juego sigue funcionando mientras la GPU trabaja (procesamiento en paralelo).
            \end{enumerate}
        \item Sin un driver adecuado, la GPU solo funcionaría en un \textbf{modo básico}.
    \end{itemize}
\end{frame}


%----------------------------------------------------------------------------------------
% EVOLUCIÓN
%----------------------------------------------------------------------------------------
\begin{frame}{Importancia y Evolución}
    \begin{itemize}
        \item Los primeros SO surgieron para automatizar y agilizar la ejecución de tareas repetitivas (procesamiento por lotes).
        \item Con el tiempo, se incorporaron funcionalidades de tiempos compartidos y entornos multiusuario.
        \item Actualmente, se espera que un SO sea robusto, seguro, flexible, multiplataforma y con interfaces intuitivas.
        \item Ejemplos de evolución:
        \begin{itemize}
            \item \textbf{UNIX} (años 70): base de la filosofía de multitarea y multiusuario.
            \item \textbf{Windows} (años 80 y 90): popularización en computadoras personales.
            \item \textbf{Linux} (años 90): distribución libre y de código abierto, muy extendida en servidores.
        \end{itemize}
    \end{itemize}
\end{frame}

%----------------------------------------------------------------------------------------
% 2. FUNCIONES DEL SISTEMA OPERATIVO
%----------------------------------------------------------------------------------------
\section{Funciones del Sistema Operativo}
\subsection{Procesos, Memoria, Dispositivos...}

\begin{frame}{Gestión de Procesos \& Sincronización}
    \begin{itemize}
        \item \textbf{Gestión de Procesos}: El SO crea, planifica y finaliza procesos, que son programas en ejecución.
        \item \textbf{Sincronización de Procesos}: Coordina la ejecución de varios procesos o hilos (threads) que comparten datos o recursos.
    \end{itemize}
\end{frame}

\begin{frame}{Gestión de Memoria}
    \begin{itemize}
        \item Asigna y libera \textbf{espacios de memoria} para procesos y datos.
        \item Implementa \textbf{memoria virtual} para simular más memoria de la que hay físicamente. (\emph{Nuetra VM utilizará la parcición de swap que creamos}
        \item Garantiza el aislamiento entre procesos, de manera que no interfieran en la memoria de otros.
    \end{itemize}
\end{frame}

\begin{frame}{Gestión de Dispositivos y Drivers}
    \begin{itemize}
        \item Un \textbf{driver} o controlador de hardware es un programa especializado que permite al SO comunicarse con un dispositivo.
        \item Simplifica y uniformiza la interacción con periféricos y hardware.
    \end{itemize}
\end{frame}

\begin{frame}{Seguridad y Control de Errores}
    \begin{itemize}
        \item Gestión de usuarios, permisos y autenticación.
        \item Manejo de excepciones, interrupciones y señales para mantener la estabilidad del sistema.
        \item Interrupción: Ctr-C en Ubuntu.
        \item Señal: SIGSEV (Segmentation Fault).
    \end{itemize}
\end{frame}

%----------------------------------------------------------------------------------------
% 3. EL KERNEL
%----------------------------------------------------------------------------------------
\section{El Kernel}
\subsection{Núcleo}

\begin{frame}{Definición de Kernel}
    \begin{itemize}
        \item Es el \textbf{núcleo} del SO y se ejecuta en modo privilegiado.
        \item Gestiona la comunicación directa con el hardware y con los procesos.
    \end{itemize}
\end{frame}

\begin{frame}{Monolítico vs. Microkernel}
    \begin{itemize}
        \item \textbf{Monolítico}: todo el SO se ejecuta en modo kernel (ej. Linux).
        \begin{itemize}
            \item Rendimiento superior.
            \item Buena integración de servicios.
            \item Menos Difícil de mantener y modificar.
        \end{itemize}
        \item \textbf{Microkernel}: sólo lo esencial en modo kernel, y el resto en espacio de usuario (ej. Minix, QNX).
        \begin{itemize}
            \item Mayor estabilidad.
            \item Ideal para sistemas críticos y embebidos.
            \item Mayor flexibibilidad.
            \item Más complejidad para sincronizar procesos.
        \end{itemize}
    \end{itemize}
\end{frame}

%----------------------------------------------------------------------------------------
% 4. CLASIFICACIÓN DE LOS SISTEMAS OPERATIVOS
%----------------------------------------------------------------------------------------
\section{Clasificación de los Sistemas Operativos}

\subsection{Por número de usuarios y tareas}
\begin{frame}{Por número de usuarios y tareas}
    \begin{itemize}
        \item \textbf{Monousuario} vs. \textbf{Multiusuario}
        \item \textbf{Monotarea} vs. \textbf{Multitarea}
    \end{itemize}
\end{frame}

\subsection{Por tipo de procesamiento (CPU)}
\begin{frame}{Por tipo de procesamiento (CPU)}
    \begin{itemize}
        \item \textbf{Monoprocesador}
        \item \textbf{Multiprocesador}
        \item \textbf{Multinúcleo}
    \end{itemize}
\end{frame}

\subsection{Sistemas Centralizados vs. Distribuidos}
\begin{frame}{Sistemas Centralizados vs. Distribuidos}
    \begin{itemize}
        \item \textbf{Centralizados}: un único equipo con todos los recursos.
        \item \textbf{Distribuidos}: varios equipos independientes que cooperan.
    \end{itemize}
\end{frame}

\begin{frame}{Sistemas en la Nube}
    \begin{itemize}
        \item Abstracción de hardware, recursos bajo demanda.
        \item Altamente escalables y con alta disponibilidad.
    \end{itemize}
    \vfill
    \centering
    Arquitectura de microservicios para servidores web.
\end{frame}

\subsection{Otras Clasificaciones}
\begin{frame}{Otras Clasificaciones}
    \begin{itemize}
        \item \textbf{Tiempo Real}
        \item \textbf{Sistemas Empotrados/Embebidos}
        \item \textbf{Sistemas Híbridos}
    \end{itemize}
\end{frame}

%----------------------------------------------------------------------------------------
% 5. EJEMPLOS REALES
%----------------------------------------------------------------------------------------
\subsection{Ejemplos Reales de Sistemas Operativos}

\begin{frame}{Windows}
    \begin{itemize}
        \item Multitarea y multiusuario.
        \item Kernel propietario bastante monolítico.
    \end{itemize}
\end{frame}

\begin{frame}{Linux}
    \begin{itemize}
        \item Kernel monolítico.
        \item Amplio soporte y aplicaciones en servidores, supercomputación y dispositivos embebidos.
    \end{itemize}
\end{frame}

\begin{frame}{macOS}
    \begin{itemize}
        \item Basado en Darwin (BSD + XNU).
        \item Integrado con hardware de Apple.
    \end{itemize}
\end{frame}

\begin{frame}{Android e iOS}
    \begin{itemize}
        \item \textbf{Android}: Kernel Linux, adaptado a dispositivos móviles.
        \item \textbf{iOS}: Derivado de Darwin (XNU), muy optimizado para hardware Apple.
    \end{itemize}
\end{frame}


%----------------------------------------------------------------------------------------
% 7. TENDENCIAS ACTUALES
%----------------------------------------------------------------------------------------
\section{Actualidad}

\subsection{Tendencias}
\begin{frame}{Virtualización y Contenedores}
    \begin{itemize}
        \item Máquinas virtuales vs. contenedores (Docker).
        \item Orquestación (Kubernetes) y escalabilidad.
    \end{itemize}
\end{frame}

\begin{frame}{Sistemas Operativos en la Nube}
    \begin{itemize}
        \item IaaS, PaaS, SaaS.
        \item Distribución de recursos y alta disponibilidad.
    \end{itemize}
\end{frame}

\subsection{Seguridad y Retos Actuales}
\begin{frame}{Seguridad y Retos Actuales}
    \begin{itemize}
        \item Cifrado, aislamiento de procesos, protección contra malware.
        \item Diseño para IoT con recursos limitados y complejidad creciente.
    \end{itemize}
\end{frame}


\end{document}
